%%%%%%%%%%%%%%%%%%%%%%%%%%%%%%%%%
%
%       EXERCÍCIO
%
%%%%%%%%%%%%%%%%%%%%%%%%%%%%%%%%%

\definevalorquestao

\begin{exercicioBanco}[\valorquestao]
Um tetraedro regular e justo (com 4 faces equiprováveis) possui suas faces pintadas com as seguintes cores:
\begin{multicols}{2}
\begin{itemize}
    \item \textbf{Face 1:} vermelho;
    \item \textbf{Face 2:} azul;
    \item \textbf{Face 3:} verde;
    \item \textbf{Face 4:} dividida em três regiões, uma vermelha, uma azul e uma verde.
\end{itemize}
\end{multicols}

\noindent O tetraedro é lançado e a face em contato com a mesa é observada. Considere os seguintes eventos:
\begin{align*}
    V &: \text{``A face sorteada contém a cor Vermelha''} \\
    A &: \text{``A face sorteada contém a cor Azul''} \\
    G &: \text{``A face sorteada contém a cor Verde''}
\end{align*}

\begin{enumerate}
    \item[(a)] Determine as probabilidades individuais de cada evento: $P(V)$, $P(A)$ e $P(G)$.
    \item[(b)] Mostre que os eventos $V$, $A$ e $G$ são independentes dois a dois.
    \item[(c)] Os eventos $V$, $A$ e $G$ são independentes? Justifique.
\end{enumerate}
\end{exercicioBanco}

